\section{Bài báo giải được gì, và để lại gì}
\label{sec:ch06-con-lai-gi}

\index{Bahdanau 2015!bài giải được gì}%
Phát biểu cho gọn: bài bỏ được ràng buộc rằng cả câu nguồn phải đi lọt qua một
vector bề rộng cố định, và bỏ bằng cách để encoder trả về một vector cho mỗi vị
trí rồi để decoder tự chọn trộn.

Con số của bài: RNNsearch-50 đạt BLEU 26.75 so với 17.82 của RNNencdec-50 trên
cùng dữ liệu và cùng quy trình huấn luyện. Trên riêng các câu không chứa từ
\tn{ngoài từ điển}{out-of-vocabulary}, 34.16 so với 26.71. Bản huấn luyện lâu
hơn, RNNsearch-50 có dấu sao, đạt 28.45 và 36.15, tức ngang ngửa hệ đối chứng
Moses ở 33.30 và 35.63 trên cột thứ hai.

Chi tiết tôi thích nhất trong bảng 1 của họ không phải dòng cao nhất.
RNNsearch-30, tức mô hình attention chỉ huấn luyện trên câu tối đa 30 từ, đạt
21.50 và vượt RNNencdec-50 ở 17.82, là mô hình vector cố định được huấn luyện
trên câu dài tới 50 từ. Kiến trúc mới với dữ liệu ngắn hơn vẫn hơn kiến trúc cũ
với dữ liệu dài hơn.

\subsection*{Cái bài tự nêu là còn thiếu}

\index{Bahdanau 2015!giới hạn bài tự nêu}%
Phần kết luận nêu đúng một chuyện: xử lý từ hiếm và từ chưa gặp. Cùng món nợ mà
chương~\ref{ch:encoder-decoder} ghi lại từ bài Sutskever, và nó còn nguyên.

Nhưng chỗ đáng đọc hơn nằm ở mục 6.1, trong phần so sánh với một bài khác, và
tôi chưa thấy ai trích nó:

\begin{quote}
\enquote{Our approach, on the other hand, requires computing the annotation
weight of every word in the source sentence for each word in the translation.
This drawback is not severe with the task of translation in which most of input
and output sentences are only 15-40 words. However, this may limit the
applicability of the proposed scheme to other tasks.}
\end{quote}

Đọc kỹ ba câu ấy. Câu thứ nhất: cơ chế cần \(T_x \times T_y\) phép chấm điểm
cho một cặp câu. Câu thứ hai: chuyện đó không nghiêm trọng, vì câu dịch chỉ dài
15 tới 40 từ. Câu thứ ba: nhưng nó có thể hạn chế việc dùng cơ chế này cho tác
vụ khác.

Bài \emph{biết} về chi phí bậc hai. Bài gọi nó là một nhược điểm, bằng chữ
\enquote{drawback} của chính họ. Rồi bài gạt nó đi bằng một lý do hoàn toàn hợp
lý vào năm 2015: ở quy mô một câu, tích ấy nhỏ.

\index{Bahdanau 2015!lý do gạt bỏ chi phí bậc hai hết đúng}%
Lý do ấy đúng cho tới khi nó hết đúng, và nó hết đúng vì hai chuyện xảy ra sau
đó. Chuyện thứ nhất là chương~\ref{ch:transformer}: Vaswani và cộng sự bỏ hẳn
mạng hồi quy và dùng attention cho \emph{mọi} chỗ, kể cả nối câu nguồn với
chính nó, nên tích \(T_x \times T_y\) thành \(n^2\) và nó xuất hiện ở mọi tầng
thay vì một lần. Chuyện thứ hai là ngữ cảnh dài ra: không còn 15 tới 40 từ mà
là hàng nghìn, rồi hàng trăm nghìn.

Chương~\ref{ch:chi-phi-bac-hai} là chương dành cho đúng câu \enquote{this may
limit the applicability of the proposed scheme to other tasks} ấy.

Tôi ghi lại chỗ này vì nó là một mẫu đáng nhớ hơn bản thân sự kiện. Bài không
bỏ sót chi phí; bài cân nhắc và kết luận đúng theo bối cảnh của nó. Cái làm
kết luận ấy hết đúng không phải một lập luận hay hơn, mà là bối cảnh đổi.

\subsection*{Cái bài không nêu, và chương này đo được}

\index{attention!chỗ attention không cứu được, nhắc lại}%
Hai chuyện, cả hai đều từ số đo của chương này chứ không từ bài.

Thứ nhất, cơ chế không làm bề rộng trạng thái hết quan trọng.
Mục~\ref{sec:ch06-be-rong} đo được rằng ở \(d_h = 16\) mô hình attention chỉ
đạt 0.0600 và ở \(d_h = 8\) thì tròn không. Cái nó gỡ bỏ là yêu cầu rằng
\emph{câu nguồn} phải vừa trong trạng thái. Trạng thái decoder vẫn phải đủ rộng
để làm việc của nó.

Thứ hai, phần đóng góp của \tn{encoder hai chiều}{bidirectional encoder} không
tách được khỏi phần đóng góp của số tham số, ít nhất là không trên corpus này.
Mục~\ref{sec:ch06-hai-chieu} đo được rằng một encoder một chiều với số tham số
tương đương cũng đạt 1.0000. Đó là phát biểu về corpus chứ không về bài, vì văn
phạm ở đây không có nhập nhằng nào để ngữ cảnh phải gỡ.

\subsection*{Cái chương~\ref{ch:lstm} để lại thì chương này vẫn không đụng tới}

\index{attention!vẫn là phụ kiện gắn lên mạng hồi quy}%
Đây là chỗ chương này giao lại cho phần sau của sách, và nó là chỗ lớn nhất.

Encoder vẫn là mạng hồi quy. Decoder vẫn là mạng hồi quy. Một câu \(T\) token vẫn là \(T\) lượt tính nối đuôi nhau, không lượt
nào bắt đầu được trước khi lượt trước xong. Encoder hai chiều còn làm nó nặng
gấp đôi ở phía nguồn.

Attention, ở bài này, là một thứ lắp thêm. Nó nằm giữa hai mạng hồi quy và làm
cho chỗ nối giữa chúng rộng ra. Nó không thay thế mạng nào cả.

Nói cách khác: chương~\ref{ch:encoder-decoder} để lại ba món nợ, chương này trả
hai. Encoder không còn vứt bỏ \(T - 1\) trạng thái, và decoder đã có cách hỏi
lại câu nguồn. Món thứ ba, tính tuần tự, còn nguyên vẹn.

\index{attention!hai năm giữa hai bài}%
Trong lịch sử thật, nó còn nguyên thêm hai năm nữa. Và cách nó được trả thì
ngược với mọi thứ chương này vừa xây: không phải bằng cách lắp thêm gì nữa, mà
bằng cách tháo mạng hồi quy ra và giữ lại đúng cái phụ kiện.

Bởi vì cái phụ kiện ấy, nhìn kỹ, không cần mạng hồi quy để chạy.
Phương trình~\eqref{eq:ch06-tong-co-trong-so} là một tổng có trọng số trên các
vị trí. Nó không có \(t-1\) ở đâu cả. Thứ duy nhất bắt nó phải chờ là các
\(\vect{h}_j\) được sinh ra tuần tự, và các \(\vect{s}_{i-1}\) cũng thế. Nếu
sinh được \(\vect{h}_j\) mà không cần đợi \(\vect{h}_{j-1}\), thì cả nửa
encoder chạy song song được ngay.

Đó là câu chương~\ref{ch:transformer} mở đầu bằng.
